\subsubsection{Summary of Discrete Object Conclusions for window-\texorpdfstring{\(6\)}{6}}\label{subsubsec:window6-discrete-results}
To highlight the ``minimal anchor point'' role of window-\(6\) and to organize the generalization to arbitrary resolution \(m\) as \emph{discrete object conclusions} (rather than interpretive narrative), several results already established in the main text and appendices can be consolidated into the following eight complementary assertions.
Together they present the closed form of structures, invariants, and discriminants as three classes of discrete objects within this framework: on one hand, the fiber histogram and Fourier spectrum of \(\Fold_m/\Fold^{\mathrm{bin}}_m\) determine statistical/information deficits; on the other hand, the boundary layer tower and Zeckendorf uniqueness compress unification dimensions into auditable resolution signatures; and furthermore, the algebraic certificates and zero-set divisor-scale counting of the collision kernel bring spectral-side constants and arithmetic modulations into a reproducible interface.

\begin{enumerate}
  \item \textbf{Stirling linear defect of the fold gauge volume: the difference between \texorpdfstring{$g_m$}{gm} and \texorpdfstring{$\kappa_m$}{kappa} tends to \(1\).}
  Let
  \(\mathcal{G}^{\mathrm{bin}}_m:=\mathrm{Aut}(\Fold^{\mathrm{bin}}_m)\)
  be the fold-gauge group induced by the binary interval fold \(\Fold^{\mathrm{bin}}_m\), and define the gauge volume density
  \[
  g_m:=\frac{1}{2^m}\log\card{\mathcal{G}^{\mathrm{bin}}_m},
  \qquad
  \kappa_m
  :=\frac{1}{2^m}\sum_{w\in X_m}d^{\mathrm{bin}}_m(w)\log d^{\mathrm{bin}}_m(w)
  =\mathbb{E}_{\pi_m}\!\bigl[\log d^{\mathrm{bin}}_m(X)\bigr],
  \quad
  \pi_m(w):=\frac{d^{\mathrm{bin}}_m(w)}{2^m}.
  \]
  Then for all \(m\ge 1\) one has the squeeze
  \(
  \kappa_m-1\le g_m\le \kappa_m
  \)
  (Proposition \ref{prop:fold-bin-gauge-volume}), and the Stirling leading term gives a linear difference expansion to exponential precision:
  \[
  g_m=\kappa_m-1+O\!\left(\frac{|X_m|\,m}{2^m}\right)
  =\kappa_m-1+O\!\left(m\left(\frac{\varphi}{2}\right)^m\right),
  \qquad
  \kappa_m-g_m=1+o(1)
  \]
  (Theorem \ref{thm:fold-bin-gauge-volume-stirling}).
  This conclusion strictly connects ``gauge volume (factorial product)'' and ``fold defect (\(d\log d\) average)'' as two normalizations of the same thermodynamic quantity, with exponentially small error.
  \par
  At the window-\(6\) anchor point, the fiber size histogram is \(2:8,3:4,4:9\) (Corollary \ref{cor:fold-bin-m6-fiber-hist}), hence
  \(
  \mathcal{G}^{\mathrm{bin}}_6\cong S_{2}^{8}\times S_{3}^{4}\times S_{4}^{9}
  \)
  and its abelianized charge lattice satisfies
  \(
  (\mathcal{G}^{\mathrm{bin}}_6)^{\mathrm{ab}}\cong (\ZZ_2)^{21}
  \)
  (Corollary \ref{cor:fold-bin-gauge-m6}).
  Through the fiber sign map \(\mathrm{sgn}_6\) (Definition \ref{def:fold-bin-gauge-sign} and Corollary \ref{cor:fold-bin-gauge-sign}), this charge lattice can be interpreted as \(21\) independent ``fiber parity charges'' (horizon bits) on \(X_6\); and the boundary sheet-flip gives a canonical \((\ZZ_2)^3\) coordinate sublattice within it (Corollary \ref{cor:fold-bin-boundary-center-z2}).

  \item \textbf{When the combinatorial upper bound for the maximum fiber saturates: explicit discriminant and finite resonance windows.}
  Write
  \(d^{\mathrm{bin}}_{\max}(m):=\max_{w\in X_m}d^{\mathrm{bin}}_m(w)\).
  The maximum is attained at \(w=0^m\), and satisfies the combinatorial upper bound
  \(d^{\mathrm{bin}}_{\max}(m)\le F_{K(m)-m+2}\);
  moreover, the bound is saturated if and only if the threshold constraint is redundant for all legal tails, i.e.,
  \[
  d^{\mathrm{bin}}_{\max}(m)=F_{K(m)-m+2}
  \quad\Longleftrightarrow\quad
  2^m-1\ge S_{\max}(m),
  \]
  where \(S_{\max}(m)\) has a completely explicit closed form (Corollary \ref{cor:fold-bin-saturation}).
  This saturation condition further forces \(2^m\) and a certain Fibonacci number to approach each other exponentially, thereby yielding the Diophantine rigidity of ``only finitely many saturations'' (Remark \ref{rem:fold-bin-finite}).
  In the full scan audit for \(m\le 400\), the set of saturation points is the finite set listed in Table \ref{tab:fold_binary_saturation_points}, forming a sequence of discrete ``resonance resolution windows.''

  \item \textbf{Fourier--energy driving mechanism for the maximum fiber and constant-level bias lower bound.}
  On the address modulo \(F_{m+2}\), let \(d_m(r)\) be the fold multiplicity function and \(\widehat d_m(k)\) its discrete Fourier coefficients,
  and write the mean \(\overline d_m=2^m/F_{m+2}\), the maximum \(M_m:=\max_r d_m(r)\), and the spectral peak \(S_m:=\max_{k\neq 0}|\widehat d_m(k)|\).
  Then the maximum fiber is quantitatively forced by the nontrivial spectral energy:
  \[
  M_m\ge \overline d_m+\frac{1}{F_{m+2}}\Biggl(\sum_{k\ne 0}\abs{\widehat d_m(k)}^2\Biggr)^{1/2}
  \quad\text{(Theorem \ref{thm:fold-max-fiber-spectrum})},
  \qquad
  M_m\ge \overline d_m+\frac{S_m}{F_{m+2}}
  \quad\text{(Theorem \ref{thm:fold-max-fiber-fourier})}.
  \]
  In particular, an auditable lower bound from a fixed Fibonacci modulus implies a constant-level multiplicative bias:
  \[
  \liminf_{m\to\infty}\frac{M_m}{\overline d_m}\ \ge\ 1+C_\varphi\ \approx\ 1.2244178274
  \qquad\text{(Corollary \ref{cor:fold-max-fiber-constant-bias})}.
  \]
  Therefore the ``excess'' of the maximum fiber is not accidental, but results from the nontrivial spectral amplitude retaining at least approximately \(22.44\%\) multiplicative bias in the limit.

  \item \textbf{Fibonacci rigidity of the top fiber spectrum: closed forms for the largest/second-largest/third-largest, proportional isolation bands, and hidden-bit splitting spectrum.}
  For \(m\ge 2\) and \(x\in X_m\), let
  \(
  d_m(x):=\card{\Fold_m^{-1}(x)}
  \), and define the top spectral levels
  \[
  \begin{aligned}
  D_m &:=\max_{x\in X_m} d_m(x),\\
  D_m^{(2)} &:=\max\{d_m(x):x\in X_m,\ d_m(x)<D_m\},\\
  D_m^{(3)} &:=\max\{d_m(x):x\in X_m,\ d_m(x)<D_m^{(2)}\}.
  \end{aligned}
  \]
  Write the set of maximum achievers \(\mathcal{M}_m:=\{x\in X_m:\ d_m(x)=D_m\}\).
  With the Fibonacci extension \(F_0=0,F_1=1\), the following hold:
  \begin{itemize}
    \item \textbf{Maximum fiber closed form, achiever construction, and two-step recurrence:} For all \(m\ge 2\),
    \[
    D_{2k}=F_{k+2},\qquad D_{2k+1}=2F_{k+1}.
    \]
    There exist explicit achiever families \(x^{(m)}\in X_m\) (writing \(m=4q+r,\ r\in\{0,1,2,3\}\)):
    \[
    x^{(m)}=
    \begin{cases}
    (0001)^{\,q-1}0000,& r=0,\ q\ge 1,\\[2pt]
    (0001)^{\,q-1}00100,& r=1,\ q\ge 1,\\[2pt]
    (0001)^{\,q}00,& r=2,\ q\ge 0,\\[2pt]
    (0001)^{\,q-1}0000100,& r=3,\ q\ge 1,
    \end{cases}
    \qquad
    d_m\!\bigl(x^{(m)}\bigr)=D_m.
    \]
    Furthermore, \(D_m\) satisfies the two-step recurrence with initial values
    \[
    D_m=D_{m-2}+D_{m-4}\quad(m\ge 6),\qquad (D_2,D_3,D_4,D_5)=(2,2,3,4)
    \]
    (Theorem \ref{thm:pom-max-fiber}; Corollary \ref{cor:pom-D-rec}).

    \item \textbf{Second-largest fiber closed form and global gap certificate (proportional isolation band):} For \(k\ge 5\),
    \[
    D_{2k}^{(2)}=F_{k+2}-F_{k-4},\qquad
    D_{2k+1}^{(2)}=2F_{k+1}-F_{k-4}
    \]
    (Theorem \ref{thm:pom-second-max-fiber-closed-form}).
    Hence for any non-maximum achiever \(x\in X_m\setminus\mathcal{M}_m\),
    \[
    d_m(x)\le D_m^{(2)}\le \varepsilon\,D_m
    \]
    with a uniform proportional isolation band, where the explicit ``one-cut'' certificate can be taken as
    \[
    \boxed{\;\frac{D_m^{(2)}}{D_m}\le \frac{25}{26}\qquad(m\ge 2)\;}
    \]
    with equality at \(m=13\); the asymptotic worst-case limiting constants are
    \[
    \lim_{k\to\infty}\frac{D_{2k+1}^{(2)}}{D_{2k+1}}
    =1-\frac{1}{11+5\sqrt5}\approx 0.9549150281,
    \qquad
    \lim_{k\to\infty}\frac{D_{2k}^{(2)}}{D_{2k}}
    =1-\frac{1}{9+4\sqrt5}\approx 0.9442719090
    \]
    (Corollary \ref{cor:pom-second-max-gap-constant}).

    \item \textbf{Third-largest fiber closed form and Fibonacci ladder gaps:} For \(k\ge 6\),
    \[
    D_{2k}^{(3)}=F_{k+2}-F_{k-3}
    \]
    (Theorem \ref{thm:pom-third-max-fiber-even-closed-form}).
    For the odd branch in the stable phase (\(m=2k+1\ge 19\)), the exact closed form is
    \[
    D_{2k+1}^{(3)}=2F_{k+1}-F_{k-4}-F_{k-8}
    \]
    (Remark \ref{rem:pom-third-max-fiber-odd-closed-form}).
    Therefore the top spectrum exhibits rigid Fibonacci ladder gaps:
    \[
    D_{2k}-D_{2k}^{(2)}=F_{k-4},\quad D_{2k}-D_{2k}^{(3)}=F_{k-3}\ (k\ge 6),
    \qquad
    D_{2k+1}-D_{2k+1}^{(2)}=F_{k-4}\ (k\ge 5),
    \]
    \[
    D_{2k+1}-D_{2k+1}^{(3)}=F_{k-4}+F_{k-8}\ (k\ge 9)
    \]
    (Remark \ref{rem:pom-top-fiber-spectrum-ladder}).

    \item \textbf{Hidden-bit splitting spectrum of maximum fiber achievers:} For any maximum achiever \(x\in X_m\) (\(d_m(x)=D_m\)), there exists a unique one-bit splitting
    \[
    d_m(x)=d_{m,0}(x)+d_{m,1}(x),
    \]
    where \(d_{m,b}(x)\) counts the two branches of preimages with hidden bit \(b\in\{0,1\}\) (Remark \ref{rem:pom-max-fiber-achievers-bsplit}; Table \ref{tab:fold_max_fiber_achievers_bsplit}).
    In the stable threshold range, the splitting vector takes only Fibonacci-type forms:
    \[
    \begin{aligned}
    m=2k:\quad &(d_{m,0},d_{m,1})\in\{(F_{k+1},F_k),(F_k,F_{k+1})\},\\
    m=2k+1:\quad &(d_{m,0},d_{m,1})\in
    \Bigl\{(2F_k,2F_{k-1}), (2F_{k-1},2F_k),\\
    &\qquad (F_{k+1},F_{k+1})\Bigr\},
    \end{aligned}
    \]
    and the balanced type \((F_{k+1},F_{k+1})\) appears with finite multiplicity as a fixed point of the exchange involution, thereby providing an auditable bisection fingerprint of the maximal phase set.
    This splitting spectrum can be viewed as the minimal sufficient statistic for the ``two-branch overflow structure'' of the top layer, and is compatible with the rigidity defect law of the top ladder gaps.

    \item \textbf{Complete structural characterization and finiteness theorem for maximum fiber achievers:}
    Under the induced graph representation of candidate sites (Theorem \ref{thm:pom-fiber-indset-factorization}), there exists an induced graph determined solely by \(x\):
    \[
    \Gamma_x\ \cong\ \bigsqcup_{i=1}^r P_{\ell_i},
    \qquad
    d_m(x)=\prod_{i=1}^r F_{\ell_i+2}.
    \]
    Hence the structure of maximum fiber achievers can be completely reduced to discriminating the component length multiset \(\{\ell_i\}\).
    In the stable threshold range (Theorem \ref{thm:pom-max-achievers-phase-stabilization}), this discrimination is further rigidified:
    \begin{itemize}
      \item \textbf{Even \(m=2k\):}\(\ x\in\mathcal{M}_m\iff \Gamma_x\cong P_k\) (single full-length path component).
      \item \textbf{Odd \(m=2k+1\):}\(\ x\in\mathcal{M}_m\iff \Gamma_x\cong P_{k-1}\sqcup P_1\) (one full-length path plus one length-\(1\) discrete component).
    \end{itemize}
    Moreover, the achiever set is closed under the reversal involution \(\mathrm{rev}\); its cardinality satisfies the uniform upper bound \(|\mathcal{M}_m|\le 4\), and in the stable phase,
    \[
    |\mathcal{M}_{2k}|=2\quad(k\ge 5),
    \qquad
    |\mathcal{M}_{2k+1}|=4\quad(k\ge 6)
    \]
    (Theorem \ref{thm:pom-max-achievers-phase-stabilization}; verifiable term by term in the auditable window of Table \ref{tab:fold_max_fiber_achievers_phase}).

    \item \textbf{Exponential amplification law with exact constants for the maximum-to-average fiber ratio:}
    Let the average fiber size be
    \[
    \overline d_m:=\frac{1}{|X_m|}\sum_{x\in X_m}d_m(x)=\frac{2^m}{|X_m|}=\frac{2^m}{F_{m+2}}.
    \]
    Then the maximum amplification ratio satisfies the limiting exponential law
    \[
    \lim_{m\to\infty}\frac{1}{m}\log\frac{D_m}{\overline d_m}
    =\log\!\left(\frac{\varphi^{3/2}}{2}\right),
    \qquad \varphi=\frac{1+\sqrt5}{2}.
    \]
    More precisely, there exist exact multiplicative constants split by parity such that
    \[
    \frac{D_m}{\overline d_m}
    =
    \begin{cases}
    \displaystyle \frac{\varphi^{4}}{5}\Bigl(\frac{\varphi^{3/2}}{2}\Bigr)^{m}\,(1+o(1)),& m\ \text{even},\\[1.2ex]
    \displaystyle \frac{2\varphi^{5/2}}{5}\Bigl(\frac{\varphi^{3/2}}{2}\Bigr)^{m}\,(1+o(1)),& m\ \text{odd},
    \end{cases}
    \]
    so the ``maximum fiber dominance'' has a strict exponential amplification rate \(\varphi^{3/2}/2>1\) with explicitly computable constant prefactors (see also Remark \ref{rem:pom-defect-gap}).

    \item \textbf{Parity criterion for fiber cardinality: \texorpdfstring{$\bmod 2$}{mod 2} determined entirely by component lengths \texorpdfstring{$\bmod 3$}{mod 3}.}
    Again let \(\Gamma_x\cong\bigsqcup_i P_{\ell_i}\). Since \(d_m(x)=\prod_i F_{\ell_i+2}\) and the Pisano period of the Fibonacci sequence modulo \(2\) is \(3\), the strict criterion is
    \[
    d_m(x)\equiv 1\pmod 2
    \quad\Longleftrightarrow\quad
    \forall i,\ \ell_i\not\equiv 1\pmod 3,
    \]
    equivalently
    \[
    d_m(x)\equiv 0\pmod 2
    \quad\Longleftrightarrow\quad
    \exists i,\ \ell_i\equiv 1\pmod 3.
    \]
    Therefore the parity of the fiber cardinality is a purely local, auditable arithmetic invariant independent of uplift/boundary certificates: it is determined solely by the \(\bmod 3\) classes of the component lengths of the invertible site graph.

    \item \textbf{Rational generating function, characteristic polynomial, and four-periodic sub-leading oscillation term for the maximum fiber sequence:}
    To close the two-step recurrence into a single unified linear recurrence, take the convention \(D_0:=1,D_1:=2\) (so that \((D_2,D_3,D_4,D_5)=(2,2,3,4)\) is consistent with Corollary \ref{cor:pom-D-rec}), then
    \[
    D_m=D_{m-2}+D_{m-4}\qquad(m\ge 4),
    \]
    with characteristic polynomial
    \[
    r^{4}-r^{2}-1=0,
    \]
    whose four characteristic roots are \(\pm\sqrt{\varphi}\) and \(\pm i/\sqrt{\varphi}\).
    The corresponding ordinary generating function is the strict rational expression
    \[
    \sum_{m\ge 0} D_m\,z^m=\frac{1+2z+z^{2}}{1-z^{2}-z^{4}}.
    \]
    Therefore the radius of convergence is \(1/\sqrt{\varphi}\) and the dominant growth rate is \(\sqrt{\varphi}\); moreover, due to the sub-leading characteristic roots \(\pm i/\sqrt{\varphi}\), the sub-leading correction term exhibits a \(4\)-periodic phase (controlled by \(i^m\)), with amplitude exactly \(O(\varphi^{-m/2})\). Equivalently, there exist constants \(A_\pm\) such that
    \[
    D_m=A_+(\sqrt{\varphi})^{m}+A_-(-\sqrt{\varphi})^{m}+O(\varphi^{-m/2}),
    \]
    thereby yielding a strict two-scale decomposition of ``dominant growth -- four-periodic fine oscillation.''

    \item \textbf{Strict sub-multiplicativity of the Fibonacci product and the ``component merging gain'' principle:}
    Fibonacci numbers satisfy the strict fusion identity and inequality (Lemma \ref{lem:pom-fib-fusion-submultiplicativity}): for any \((a,b)\neq(1,1)\),
    \(
    F_aF_b< F_{a+b-1}<F_{a+b}
    \).
    Hence in the component multiplicative decomposition \(d_m(x)=\prod_i F_{\ell_i+2}\) caliber, for any two component lengths \(\ell,\ell'\ge 1\), if they can be joined on the candidate chain by filling in one bridging site into a longer component, the fiber cardinality exhibits a strict ``merging gain''
    \[
    F_{\ell+2}F_{\ell'+2}<F_{\ell+\ell'+3},\qquad
    F_{\ell+\ell'+3}-F_{\ell+2}F_{\ell'+2}=F_{\ell+1}F_{\ell'+1}\ge F_{\min\{\ell,\ell'\}+1}
    \]
    (Corollary \ref{cor:pom-fib-component-fusion-gain}).
    This gain lower bound provides a purely combinatorial, auditable stability certificate: the magnitude of the fiber cardinality deviation from saturation can be jointly controlled by the ``number of splits'' and the shortest component length.

    \item \textbf{Sub-exponential complexity of distinct fiber cardinality values: a value-range contrast at the \texorpdfstring{\(\exp(\Theta(\sqrt m))\)}{exp(Theta(sqrt m))} level.}
    Let \(\mathcal{V}_m:=\{d_m(x):x\in X_m\}\) be the value set of fiber cardinalities at resolution \(m\).
    Under the component decomposition caliber, \(|\mathcal{V}_m|\) is squeezed by integer partition numbers at the exponential scale, hence satisfies the strict sub-exponential scaling
    \[
    |\mathcal{V}_m|=\exp(\Theta(\sqrt m)),
    \]
    with explicitly available partition number upper and lower bounds (see Remark \ref{rem:pom-fiber-value-set-complexity}).
    Therefore the following structural contrast emerges: \(|X_m|=F_{m+2}=\exp(\Theta(m))\) is of exponential size, while \(\mathcal{V}_m\) is only of sub-exponential size.

    \item \textbf{Extremal structure of the maximal chain count and ``the opposite extreme from the maximum fiber'': all isolated components are optimal.}
    In the distributive lattice model \(\mathcal L_x\cong J(\bigsqcup_i Z_{\ell_i})\), the maximal chain count (equivalently, the number of linear extensions) has the closed form
    \[
    |\mathrm{MC}(\mathcal L_x)|
    =\binom{N}{\ell_1,\dots,\ell_r}\prod_{i=1}^r E_{\ell_i}
    =N!\prod_{i=1}^r\frac{E_{\ell_i}}{\ell_i!},
    \qquad N=\sum_{i=1}^r \ell_i
    \]
    (Theorem \ref{thm:pom-fence-maxchains-euler}).
    Hence under the constraint of fixed total length \(N\),
    \[
    |\mathrm{MC}(\mathcal L_x)|\le N!,
    \]
    with equality if and only if \(\ell_1=\cdots=\ell_r=1\) (all components isolated) (Corollary \ref{cor:pom-fence-maxchains-all-isolated-optimal}).
    This extremal structure is complementary to the ``single long component'' structure of the maximum fiber achievers (the rigidity characterization of \(\Gamma_x\) above): the fiber simultaneously carries the separation of ``volume maximum'' and ``scheduling freedom maximum.''

    \item \textbf{Convex hull theorem for the dual-invariant reachable domain: the Pareto front of \texorpdfstring{\((\log d,\log\mathrm{MC})\)}{(log d, log MC)} is generated by finitely many atomic lengths.}
    Under the above two multiplicative decompositions, for the normalized dual invariants
    \[
    A=\frac{1}{N}\log d_m(x),\qquad
    B=\frac{1}{N}\bigl(\log|\mathrm{MC}(\mathcal L_x)|-\log(N!)\bigr)
    \]
    (\(N=\sum_i\ell_i\)), their scaling limit as \(N\to\infty\) forms a compact convex reachable domain, and this domain equals the closed convex hull of the atomic point set
    \(
    \{(\ell^{-1}\log F_{\ell+2},\,\ell^{-1}\log(E_\ell/\ell!)):\ell\in\NN\}
    \)
    (Theorem \ref{thm:pom-fiber-ab-convex-hull}).
    Hence the Pareto front can be fully characterized by supporting hyperplanes and is generated by mixtures of finitely many atomic lengths.

    \item \textbf{Spectral rigidity of the coefficient \(4\): half-odd-integer angular quantization of the discrete Green kernel.}
    In the fence-order polynomial kernel operator caliber, the Green kernel
    \(K_k(i,j)=\min(i,j)\) has the spectral closed form
    \(
    \lambda_p(k)=\bigl(2(1-\cos\theta_p)\bigr)^{-1}=(4\sin^2(\theta_p/2))^{-1}
    \)
    with half-odd-integer angular quantization
    \(
    \theta_p=\frac{(2p-1)\pi}{2k+1}
    \)
    (Lemma \ref{lem:pom-Kk-eigenvalues}; Remark \ref{rem:pom-Kk-factor4-spectral-rigidity}).
    Therefore the coefficient \(4\) and the scale \(4k+2\) in the spectral representation are not arbitrary normalizations but are structurally forced by the mixed boundary quantization and the half-angle identity.

    \item \textbf{``Quarter-turn'' nearest singularity law: the \texorpdfstring{\(2/\pi\)}{2/pi} constant of scheduling entropy and the analytic origin of the coefficient \(4\).}
    The exponential constant of the Euler zigzag numbers \(E_n\) is
    \(
    \lim_{n\to\infty}\frac1n\log(E_n/n!)=\log(2/\pi)
    \),
    dominated by the nearest singularities of \(\sec z+\tan z\) at \(z=\pm\pi/2\) (Theorem \ref{thm:pom-fence-volume-log2pi}).
    Since the fundamental period is \(2\pi\), \(\pi/2=(2\pi)/4\) gives the strict ``quarter-turn'' scale,
    thereby unifying \(\log(2/\pi)=-\log(\pi/2)\) with the spectral-side coefficient \(4\) under the same angular scale ratio (Remark \ref{rem:pom-quarter-turn-pi2-scale}).

    \item \textbf{Homologous closure of two entropy constants: \texorpdfstring{\(\pi/2\)}{pi/2} as the universal scale of fold thermodynamics.}
    The multi-chain refinement spectral density satisfies the closed form
    \[
    h_{\mathrm{ref}}(k)=-\log\!\Bigl(2\sin\frac{\pi}{4k+2}\Bigr)
    =\log(2k+1)-\log\pi+\frac{\pi^2}{24(2k+1)^2}+\frac{\pi^4}{2880(2k+1)^4}+O(k^{-6})
    \]
    (Corollary \ref{cor:pom-fence-multichain-entropy-pi}; Remark \ref{rem:pom-quarter-turn-pi2-scale}).
    Juxtaposing with the scheduling/volume-side limiting density \(\log(2/\pi)=-\log(\pi/2)\) reveals: \(-\log\pi\) and \(\log(2/\pi)\) share the unique angular scale \(\pi/2\), realized respectively through half-angle quantization and nearest singularity dominance.

    \item \textbf{Spectral-geometric homology paradigm: endpoint regularity \texorpdfstring{\(\Longleftrightarrow\)}{<->} half-odd-integer modes \texorpdfstring{\(\Longleftrightarrow\)}{<->} quarter-turn constants.}
    In summary, the regularity condition at mixed boundary endpoints forces half-odd-integer modes and introduces the half-angle coefficient \(4\) on the discrete spectral side, while on the analytic combinatorics side it locks \(\log(2/\pi)\) via the \(\pi/2\) quarter-turn nearest singularity;
    the two jointly constitute an implementation-independent ``half-angle/quarter-turn'' homology mechanism, thereby providing a unified mathematical prototype for the \(1/4\)-type proportional constants that recurrently appear in the horizon endpoint-related passages of this paper (e.g., the \(1/4\) form in the black hole horizon semantics) (see the attribution paragraph of Remark \ref{rem:pom-quarter-turn-pi2-scale}).

    \item \textbf{Determinantal rigidity of the Gram decomposition and the \texorpdfstring{\((4\sin^2)\)}{(4 sin^2)} product identity.}
    The Green kernel \(K_k(i,j)=\min(i,j)\) admits the integer Gram decomposition \(K_k=B_kB_k^{\mathsf T}\) (\(B_k\) being the unit lower triangular all-\(1\) matrix), hence \(\det(K_k)=1\) (Lemma \ref{lem:pom-Kk-gram-det}).
    Combined with the eigenvalue closed form \(\lambda_p=(4\sin^2\phi_p)^{-1}\) (\(\phi_p=(2p-1)\pi/(4k+2)\)), one obtains the strict product identity
    \(
    \prod_{p=1}^{k}4\sin^2(\phi_p)=1
    \), and from \(\mathrm{tr}(K_k)=k(k+1)/2\) the corresponding strict summation identity (Corollary \ref{cor:pom-Kk-sine-product-sum}).

    \item \textbf{Chebyshev identification: \texorpdfstring{\(\lambda/4\)}{lambda/4} as the unique integer-coefficient normalization scale.}
    Write \(L_k:=K_k^{-1}\) for the mixed-boundary discrete Laplacian, then its characteristic polynomial can be written as
    \[
    \det(\lambda I-L_k)=(-1)^k\,P_k\!\Bigl(1-\frac{\lambda}{4}\Bigr),
    \qquad
    P_k(y):=\frac{T_{2k+1}(\sqrt y)}{\sqrt y}\in\ZZ[y],
    \]
    so the spectral points satisfy \(\mu_p=4\sin^2(\phi_p)\) (Theorem \ref{thm:pom-Lk-chebyshev-charpoly}).
    This representation promotes the coefficient \(4\) to a joint constraint of ``Chebyshev modal quantization + integer-coefficient spectral algebra'': \(\lambda\mapsto 1-\lambda/4\) is the natural linear normalization that maps spectral points to \(\cos^2(\phi_p)\) and thereby aligns with \(T_{2k+1}\).

    \item \textbf{Closed form for the shifted determinant and strict linear free energy (area law form).}
    For \(t>-4\), the strict identity holds:
    \[
    \det(L_k+tI)=\det(I+tK_k)=P_k\!\Bigl(1+\frac{t}{4}\Bigr)=\frac{T_{2k+1}(\sqrt{1+t/4})}{\sqrt{1+t/4}},
    \]
    yielding the free energy density limit
    \[
    \lim_{k\to\infty}\frac{1}{2k+1}\log\det(L_k+tI)=\mathrm{arcosh}\!\Bigl(\sqrt{1+\frac{t}{4}}\Bigr),
    \]
    and the explicit renormalization constant term
    \(
    \lim_{k\to\infty}F_k(t)=-\log\!\bigl(2\sqrt{1+t/4}\bigr)
    \)
    (Corollary \ref{cor:pom-Lk-shifted-det-free-energy}).

    \item \textbf{Fully explicit spectral gap constants: the rigid appearance of \texorpdfstring{\(\pi^2/(2k+1)^2\)}{pi^2/(2k+1)^2}.}
    The smallest/second-smallest eigenvalues of \(L_k\) satisfy
    \[
    \mu_1=4\sin^2\!\Bigl(\frac{\pi}{4k+2}\Bigr),\qquad
    \mu_2=4\sin^2\!\Bigl(\frac{3\pi}{4k+2}\Bigr),
    \]
    hence
    \(
    \mu_1=\pi^2/(2k+1)^2+O(k^{-4})
    \)
    and \(\lambda_{\max}(K_k)=\mu_1^{-1}=(2k+1)^2/\pi^2+O(1)\)
    (Corollary \ref{cor:pom-Lk-gap-pi2}).

    \item \textbf{Bernoulli--polynomial law for higher moments: all odd-angle \texorpdfstring{\(\csc^{2r}\)}{csc^{2r}} power sums are even-degree polynomials in \texorpdfstring{\((2k+1)\)}{(2k+1)}.}
    For \(r\ge 1\) define
    \(
    S_r(k)=\sum_{p=1}^{k}\csc^{2r}(\phi_p)
    \)
    and \(n:=2k+1\), then
    \[
    S_r(k)=-\frac12+\sum_{j=1}^{r} a_{r,j}\,n^{2j},
    \qquad
    a_{r,r}=\frac{2^{2r-1}(2^{2r}-1)\,|B_{2r}|}{(2r)!},
    \]
    and \(\mathrm{tr}(K_k^{\,r})=\frac{1}{4^{\,r}}S_r(k)\), so the coefficient \(4\) systematically appears as \(4^{-r}\) across all trace moments (Proposition \ref{prop:pom-csc-moments-bernoulli}; see also Remark \ref{rem:pom-csc-moments-examples} for the first three orders in closed form).

    \item \textbf{Limiting law of the spectral measure: the empirical spectral distribution converges to the arcsine law on \texorpdfstring{\([0,4]\)}{[0,4]}.}
    The empirical spectral measure of the mixed-boundary discrete Laplacian \(L_k=K_k^{-1}\),
    \(
    \nu_k=\frac1k\sum_{p=1}^k\delta_{\mu_p}
    \),
    converges weakly to the arcsine distribution on \([0,4]\): \(\dd\nu(\mu)=\frac{1}{\pi\sqrt{\mu(4-\mu)}}\dd\mu\), yielding the trace functional limit
    \(
    \frac1k\tr f(L_k)\to \frac1\pi\int_0^4 f(\mu)/\sqrt{\mu(4-\mu)}\,\dd\mu
    \)
    (Theorem \ref{thm:pom-Lk-arcsine-law}).
    This limiting law promotes the coefficient \(4\) to a rigid spectral support endpoint scale.

    \item \textbf{Exact remainder bound for the discrete Weyl law: half-odd-integer grid error compressed to constant order.}
    For the counting function \(N_k(\mu)=\#\{p\le k:\mu_p\le\mu\}\) and
    \(
    \theta(\mu)=\arccos(1-\mu/2)
    \),
    \[
    N_k(\mu)=\left\lfloor \frac{2k+1}{2\pi}\theta(\mu)+\frac12\right\rfloor,\qquad
    \Bigl|N_k(\mu)-\frac{2k+1}{2\pi}\theta(\mu)\Bigr|\le 1
    \]
    (Corollary \ref{cor:pom-Lk-discrete-weyl}), so the empirical spectral distribution function approximates the continuous arcsine law uniformly at rate \(O(1/k)\) over the entire interval.

    \item \textbf{Dirichlet factor of the spectral \texorpdfstring{\(\zeta\)}{zeta} leading term: analytic fingerprint of the odd-mode projection.}
    For the spectral \(\zeta\) function \(\zeta_k(s)=\sum_p\mu_p^{-s}\) with \(\Re(s)>\tfrac12\), the dominant asymptotic is
    \[
    \zeta_k(s)=\Bigl(\frac{2k+1}{\pi}\Bigr)^{2s}\Bigl(1-2^{-2s}\Bigr)\zeta(2s)+O(k),
    \]
    where the factor \(\bigl(1-2^{-2s}\bigr)\) is forcibly produced by the half-odd-integer modes (odd frequency sub-spectrum) (Proposition \ref{prop:pom-Lk-spectral-zeta-dirichlet}).

    \item \textbf{Closed-form limit of the heat kernel density: Bessel--arcsine duality and diffusion scale normalization.}
    The heat trace \(H_k(t)=\tr(e^{-tL_k})\) satisfies
    \[
    \frac1k H_k(t)\ \longrightarrow\ \frac{1}{\pi}\int_0^\pi e^{-2t(1-\cos\theta)}\,\dd\theta
    =e^{-2t}I_0(2t),
    \qquad
    \frac{1}{2k+1}H_k(t)\ \longrightarrow\ \frac12 e^{-2t}I_0(2t)
    \]
    (Corollary \ref{cor:pom-Lk-heat-kernel-bessel}), so the arcsine spectral law under Laplace transform yields a strict Bessel closed form.

    \item \textbf{Unique scale of the discrete--continuous spectral isomorphism: ND determinant limit \texorpdfstring{\(\cosh(\sqrt{s})\)}{cosh(sqrt(s))}.}
    The continuous ND operator \(\mathcal L=-\dd^2/\dd x^2\) (\(u'(0)=0,u(1)=0\)) satisfies \(\det(\mathcal L+s)=\cosh(\sqrt{s})\).
    The discrete operator under the scaling
    \(
    \widetilde L_k=\frac{(2k+1)^2}{4}L_k
    \)
    achieves strict spectral matching, with the determinant limit
    \[
    \det\!\Bigl(L_k+\frac{4s}{(2k+1)^2}I\Bigr)\ \longrightarrow\ \cosh(\sqrt{s})\qquad(s\ge 0),
    \]
    so the scale factor \(t/4\) is structurally forced by the half-odd-integer mode matching of the continuous ND spectrum (Corollary \ref{cor:pom-Lk-continuum-ND-det-cosh}).

        \item \textbf{Deterministic ``area law'' decomposition: \texorpdfstring{\(\det(L_k+tI)\)}{det(L_k+tI)} in \(q\)-form with exponentially small finite-size corrections.}
    Take \(t>0\), let
    \(
    \eta(t)=\mathrm{arsinh}\sqrt{t/4}
    \)
    and \(q(t)=e^{-2\eta(t)}\in(0,1)\).
    Then the following strict identity holds:
    \[
    \det(L_k+tI)=\frac{\cosh((2k+1)\eta(t))}{\cosh(\eta(t))}
    =q(t)^{-k}\frac{1+q(t)^{2k+1}}{1+q(t)},
    \]
    so that
    \(
    \log\det(L_k+tI)=2k\eta(t)-\log(1+q(t))+\log(1+q(t)^{2k+1})
    \)
    , with the remainder satisfying \(0<\log(1+q^{2k+1})<q^{2k+1}\) (Corollary \ref{cor:pom-Lk-area-law-qform}).

    \item \textbf{Strong Szeg\H{o}-type bulk--surface decomposition with integrable closed form: the integral density of the symbol \texorpdfstring{\(a_t(\theta)\)}{a_t(theta)} is exactly \texorpdfstring{\(2\eta(t)\)}{2 eta(t)}.}
    Write
    \(a_t(\theta)=1+\frac{t}{4\sin^2(\theta/2)}\) (\(\theta\in(0,\pi)\)).
    From \(\det(K_k)=1\), one obtains \(\det(I+tK_k)=\det(L_k+tI)\), with the strict decomposition
    \[
    \log\det(I+tK_k)
    =k\cdot \frac{1}{\pi}\int_{0}^{\pi}\log a_t(\theta)\,\dd\theta
    +\log\frac{1+q(t)^{2k+1}}{1+q(t)},
    \]
    where the bulk integral term can be made explicit as
    \(
    \frac{1}{\pi}\int_{0}^{\pi}\log\!\bigl(1+\frac{t}{4\sin^2(\theta/2)}\bigr)\,\dd\theta
    =2\eta(t)
    \)
    (Proposition \ref{prop:pom-Kk-strong-szego-bulk-surface}).

    \item \textbf{Analytically auditable resolvent trace: hyperbolic closed form of \texorpdfstring{\(\tr(L_k+tI)^{-1}\)}{tr((L_k+tI)^{-1})}.}
    The resolvent trace satisfies the strict identity
    \[
    \tr(L_k+tI)^{-1}
    =\frac{2k+1}{2\sqrt{t(4+t)}}\,\tanh\bigl((2k+1)\eta(t)\bigr)-\frac{1}{2(4+t)},
    \]
    hence the trace density limit is
    \(
    \lim_{k\to\infty}\frac{1}{2k+1}\tr(L_k+tI)^{-1}=\frac{1}{2\sqrt{t(4+t)}}
    \)
    (Corollary \ref{cor:pom-Lk-resolvent-trace-hyperbolic}).

    \item \textbf{Complete closed form of Green kernel matrix entries and two-point propagator compression: hyperbolic product formula for \texorpdfstring{\((L_k+tI)^{-1}_{ij}\)}{(L_k+tI)^{-1}_{ij}}.}
    For \(1\le i\le j\le k\) the strict closed form is
    \[
    (L_k+tI)^{-1}_{ij}
    =\frac{\sinh(2i\eta(t))\,\cosh((2(k-j)+1)\eta(t))}
    {\sinh(2\eta(t))\,\cosh((2k+1)\eta(t))},
    \qquad (L_k+tI)^{-1}_{ji}=(L_k+tI)^{-1}_{ij}.
    \]
    In particular,
    \(
    (L_k+tI)^{-1}_{11}=\cosh((2k-1)\eta)/\cosh((2k+1)\eta)
    \)
    and
    \(
    (L_k+tI)^{-1}_{1k}=\cosh(\eta)/\cosh((2k+1)\eta)
    \)
    are fully explicit (Proposition \ref{prop:pom-Lk-green-kernel-entries}).

    \item \textbf{Exponential clustering and correlation length: strict ``gap law'' of the mass parameter \texorpdfstring{\(t\)}{t}.}
    For any \(1\le i\le j\le k\) the uniform upper bound holds:
    \[
    (L_k+tI)^{-1}_{ij}\le \frac{1}{\sinh(2\eta(t))}\,e^{-2\eta(t)(j-i)},
    \]
    so the correlation length can be taken as
    \(
    \xi(t)=1/(2\eta(t))=1/(2\,\mathrm{arsinh}\sqrt{t/4})
    \)
    (Corollary \ref{cor:pom-Lk-exponential-clustering-gaplaw}).

    \item \textbf{Boundary fixed point in the semi-infinite limit: algebraic closed form and Riccati equation for the reflection coefficient \texorpdfstring{\(q(t)\)}{q(t)}.}
    The boundary Green value satisfies
    \(
    \lim_{k\to\infty}(L_k+tI)^{-1}_{11}=q(t)=e^{-2\eta(t)}
    \)
    , and
    \[
    q(t)=\Bigl(\frac{\sqrt{4+t}-\sqrt{t}}{2}\Bigr)^2
    =1+\frac{t}{2}-\frac12\sqrt{t(4+t)},
    \qquad
    q^2-(2+t)q+1=0,
    \]
    thereby forming an algebraic certificate for the semi-infinite continued-fraction/Riccati problem (Corollary \ref{cor:pom-Lk-boundary-fixed-point-riccati}).

    \item \textbf{Surface free energy and boundary response: closed-form constant term and observable derivative.}
    Let
    \(f(t)=\lim_{k\to\infty}\frac1k\log\det(L_k+tI)=2\eta(t)\)
    and
    \(g(t)=\lim_{k\to\infty}(\log\det-k f(t))=-\log(1+q(t))\), then the exact decomposition reads
    \[
    \log\det(L_k+tI)=k f(t)+g(t)+\log(1+q(t)^{2k+1}),
    \]
    with
    \(
    g'(t)=\frac{q(t)}{1+q(t)}\cdot \frac{1}{\sqrt{t(4+t)}}
    \)
    (Corollary \ref{cor:pom-Lk-surface-free-energy}; see also Remark \ref{rem:pom-fence-green-kernel-area-law-audit} and the finite-window audit Table \ref{tab:fence_green_kernel_area_law_resolvent_audit}).



        \item \textbf{``Golden coupling'' uniqueness: equivalent characterizations of \texorpdfstring{\(t=1\)}{t=1} and the golden locking of the spectral free energy density.}
    Define \(\eta(t)=\mathrm{arsinh}\sqrt{t/4}\) (\(t>0\)), then the conditions
    \(
    \eta(t)=\log\varphi
    \)
    and
    \(
    q(t)=e^{-2\eta(t)}=\varphi^{-2}
    \)
    and
    \(
    \lim_{k\to\infty}\frac1k\log\det(L_k+tI)=2\log\varphi
    \)
    are strictly equivalent and yield the same unique positive solution \(t_\star=1\) (Corollary \ref{cor:pom-Lk-golden-coupling-unique}).

    \item \textbf{Fibonacci determinant and fully explicit Green kernel at the golden coupling: unification by the denominator \texorpdfstring{\(F_{2k+1}\)}{F_{2k+1}}.}
    Let \(A_k=L_k+I\), then
    \(
    \det(A_k)=F_{2k+1}
    \)
    and for \(1\le i\le j\le k\),
    \[
    (A_k^{-1})_{ij}=\frac{F_{2i}\,F_{2(k-j)+1}}{F_{2k+1}},
    \qquad (A_k^{-1})_{ji}=(A_k^{-1})_{ij},
    \]
    so the two-point correlation falls entirely into the rational number field with a single denominator \(F_{2k+1}\) (Corollary \ref{cor:pom-Lk-t1-fibonacci-det-green}).

    \item \textbf{Exact realization of ``area term + constant term + exponentially small correction'' on Fibonacci numbers.}
    By the Binet formula in its exact form for odd indices,
    \[
    \log F_{2k+1}=(2k+1)\log\varphi-\frac12\log 5+\log(1+\varphi^{-4k-2}),
    \qquad 0<\log(1+\varphi^{-4k-2})<\varphi^{-4k-2},
    \]
    so this decomposition is a strict identity at the object level rather than an asymptotic approximation (Remark \ref{rem:pom-Lk-fibonacci-area-law-exact}).

    \item \textbf{Homologous closure of boundary entropy response and Parry mass: \texorpdfstring{\(g'(1)=\pi(1)/\sqrt5\)}{g'(1)=pi(1)/sqrt5}.}
    The Parry mass of the golden mean shift satisfies
    \(
    \pi(1)=1/(\varphi^2+1)
    \)
    (Remark \ref{prop:parry-measure-golden}).
    On the spectral side, the derivative of the surface free energy \(g(t)=-\log(1+q(t))\) at \(t=1\) satisfies
    \(
    g'(1)=\pi(1)/\sqrt5
    \)
    (Corollary \ref{cor:pom-Lk-t1-parry-surface-derivative}).

    \item \textbf{Limiting distribution of Fisher zeros: arcsine law on \texorpdfstring{\([-4,0]\)}{[-4,0]} and closed-form external logarithmic potential.}
    The zeros of the polynomial \(P_k(t)=\det(L_k+tI)\) are
    \(
    t_p=-4\sin^2(\frac{(2p-1)\pi}{4k+2})\in(-4,0)
    \);
    the empirical measure of the zeros converges weakly to the arcsine distribution on \([-4,0]\):
    \(
    \dd\omega(t)=\frac{1}{\pi\sqrt{(-t)(4+t)}}\,\dd t
    \)
    , and for \(t\in\CC\setminus[-4,0]\) the external logarithmic potential has the closed form
    \(
    \lim_{k\to\infty}\frac1k\log|\det(L_k+tI)|=2\Re\,\mathrm{arsinh}\sqrt{t/4}
    \)
    (Theorem \ref{thm:pom-Lk-fisher-zeros-arcsine}).

    \item \textbf{Conformal uniformization of zeros: ``odd-root unification'' under the Joukowsky map.}
    Introducing \(w+w^{-1}=t+2\) (i.e., \(t=w+w^{-1}-2\)), the equation \(P_k(t)=0\) is equivalent to \(w^{2k+1}=-1\);
    this identifies \([-4,0]\) as the exact image domain endpoints of the unit circle under the Joukowsky map, and interprets the arcsine law as the pushforward limit of uniformly distributed discrete points on the unit circle (Remark \ref{rem:pom-Lk-fisher-zeros-uniformization}).

    \item \textbf{Riccati recursion for the boundary reflection coefficient and Fibonacci convergence rate at the golden coupling.}
    The boundary Green value \(q_k(t)=(L_k+tI)^{-1}_{11}\) satisfies the recursion
    \(q_{k+1}(t)=1/(t+2-q_k(t))\), with exact error
    \(
    0<q_k(t)-q(t)=\frac{(1-q(t)^2)q(t)^{2k}}{1+q(t)^{2k+1}}=\Theta(q(t)^{2k})
    \).
    At \(t=1\),
    \(q_k(1)=F_{2k-1}/F_{2k+1}\to \varphi^{-2}\)
    with error of order \(\Theta(\varphi^{-4k})\) (Proposition \ref{prop:pom-Lk-boundary-riccati-recursion}; see also Remark \ref{rem:pom-fence-green-kernel-golden-coupling-audit} and the audit Table \ref{tab:fence_green_kernel_golden_coupling_audit}).

  \end{itemize}

  \item \textbf{Uniqueness of the Zeckendorf--Fibonacci resolution signature \texorpdfstring{$\Sigma(\mathfrak g)$}{Sigma(g)} and the minimal solution of the NAP selector.}
  For a candidate unification Lie algebra \(\mathfrak g\), let \(D=\dim\mathfrak g\), and write the Zeckendorf unique decomposition as
  \(D=\sum_{k\in S}F_k\) (\(S\) contains no adjacent indices).
  Define the resolution signature
  \[
  \Sigma(\mathfrak g):=\{\,m=k+2:\ k\in S\,\}.
  \]
  Since the boundary layer counts satisfy \(\card{X_m^{\mathrm{bdry}}}=F_{m-2}\), one has
  \(D=\sum_{m\in\Sigma(\mathfrak g)}\card{X_m^{\mathrm{bdry}}}\), and by Zeckendorf uniqueness, \(\Sigma(\mathfrak g)\) is also unique (see \S\ref{subsec:bdry-tower-zeck-gut} and Table \ref{tab:zeckendorf_gut_signatures}).
  On this basis, imposing the non-abelian preservation constraint (NAP) \(\{6,8\}\subset \Sigma(\mathfrak g)\) and minimizing \(D\) over the simple Lie algebra families yields the minimal solution
  \(
  \mathfrak g=\mathfrak{so}(10),\ D=45,\ \Sigma(\mathfrak g)=\{6,8,11\}
  \)
  (Table \ref{tab:nap_selector_so10} and Equation \ref{eq_nap_selector_so10}).
  This construction promotes ``unification dimension'' to an auditable discrete layer signature functor and yields a fully discrete, reproducible unification algebra selector.

  \item \textbf{Multifractal skeleton of the high-degeneracy tail: isomorphic interfaces of the upper envelope (Legendre) and lower envelope (Paley--Zygmund).}
  Let \(\gamma_m(x)=(1/m)\log d_m(x)\), and define the spectral function \(f(\gamma)\) as in Definition \ref{def:pom-fiber-spectrum}.
  Then the moment pressure \(\tau(q)\) (or in the verified caliber, \(\tau(q)=\log r_q\)) gives the strict upper envelope
  \[
  f(\gamma)\ \le\ \inf_{q\ge 2}\bigl(\tau(q)-q\gamma\bigr)
  \qquad\text{(Proposition \ref{prop:pom-spectrum-legendre-upper})},
  \]
  which under standard thermodynamic formalism conditions can be promoted to equality, thereby compressing the spectral shape into computable information from \(\{r_q\}\).
  Simultaneously, the Paley--Zygmund lower envelope is directly given by \((r_q,r_{2q})\) (Proposition \ref{prop:pom-spectrum-pz-lower}), with the deficit
  \(
  \delta(q):=\log\varphi-\underline f(q)
  \)
  quantifying the decay of the ``fat-tail guarantee exponent'' (Corollary \ref{cor:pom-spectrum-intermittency-deficit}).
  At the audit points \(q=2,\dots,10\), \(\delta(q)\) rises significantly with increasing \(q\); typical values are given in Table \ref{tab:fold_collision_multifractal_envelope_pz_q10} and Figure \ref{fig:fold-collision-multifractal-envelope-pz} (Remark \ref{rem:pom-spectrum-pz-numerics}).
  Moreover, from the low-order moments \((\mu_m,\sigma_m^2)\) at finite resolution, one obtains a ``typical window plateau'' count lower bound
  \[
  \card{A_m(\varepsilon)}
  \ \ge\
  \left(1-\frac{\sigma_m^2}{m\varepsilon^2}\right)
  \exp\!\bigl(m(\log 2-\mu_m-\varepsilon)\bigr),
  \qquad
  A_m(\varepsilon):=\{x\in X_m:\ \abs{\gamma_m(x)-\mu_m}\le \varepsilon\}
  \]
  (Proposition \ref{prop:pom-spectrum-finite-m-plateau}), providing a fully auditable finite-\(m\) interface for the juxtaposition of fat tails and typical sets.

  \item \textbf{Algebraic certificate for the main-pole residue constant of the collision kernel: determined solely by \texorpdfstring{$(\chi_A,r)$}{(chiA,r)}.}
  For the adjacency matrix \(A\) of a finite-state collision kernel, let \(\chi_A(\lambda)=\det(\lambda I-A)\) (monic), and suppose the Perron root \(r=\rho(A)\) is a simple root.
  Define
  \(
  \zeta_A(z)=\det(I-zA)^{-1}
  \)
  and the main-pole residue constant
  \[
  C(A):=\lim_{z\to 1/r}(1-rz)\,\zeta_A(z).
  \]
  Then the closed-form identity holds:
  \[
  C(A)=\frac{r^{d-1}}{\chi_A'(r)}
  =\prod_{\lambda\in\sigma(A)\setminus\{r\}}\left(1-\frac{\lambda}{r}\right)^{-1},
  \qquad d:=\deg\chi_A,
  \]
  so once \(\chi_A\) and \(r\) are independently audited, \(C(A)\) can be exactly reconstructed without computing the full spectrum; furthermore, one can provide an integer-coefficient minimal polynomial-level certificate (see Definition \ref{def:pom-collision-zeta-a4} and the audit outputs in Tables \ref{tab:collision_kernel_A2_finite_part}--\ref{tab:collision_kernel_A4_finite_part}).

  \item \textbf{Exact ``divisor-scale'' counting of the zero set \texorpdfstring{$\card{\mathcal{Z}_m}$}{|Z_m|}: from scanning frequencies to scanning divisors of \texorpdfstring{$m+2$}{m+2}.}
  When \(F_{m+2}\) is even, the zero set \(\mathcal{Z}_m\) can be expressed as a union of arithmetic residue classes, so its exact computation can be reduced in dimension to:
  enumerating the few divisors \(d\mid(m+2)\) and performing finitely many divisibility screenings, constructing the corresponding congruence classes and counting with finite intersection compatibility criteria, and finally completing the union count via inclusion--exclusion.
  This procedure reduces ``scanning all frequencies \(k\)'' to ``scanning the divisor structure of \(m+2\),'' and can be directly embedded into the audit pipeline (see Remark \ref{rem:fold-zero-coset-algorithm}; implementation script \texttt{\detokenize{scripts/exp_fold_zero_coset_union_count.py}}).
\end{enumerate}

